\section*{\texorpdfstring{\S\CurrentExerciseGroup}{第{\CurrentExerciseGroup}节}}
\phantomsection
\addcontentsline{toc}{section}{第{\CurrentExerciseGroup}节习题}

\begin{enumerate}
\def\labelenumi{\arabic{enumi})}
\item
  设 \(\mathbf{T}\) 与 \(\mathbf{X}\) 为两个局部紧空间，\(\mu\) 为一个测度 \(\geqslant 0\)，定义在 \(\mathbf{T}\) 上，\(\Lambda : t \mapsto \lambda_t\) 为一个 \(\mu\) -适宜映射，把 \(\mathbf{T}\) 映入 \(\mathcal{M}_+(X)\)，\(g\) 为关于测度 \(\nu = \int \lambda_t d\mu(t)\) 局部可积的函数。试给出一个例子说明：若 \(\mathbf{X}\) 不是无穷远处可数的，则 \(g\) 可能不是局部 \(\lambda_t\) -可积的，无论 \(t\) 取何值（参见 §3, 习题 4）。
\item
  设 \(g\) 为 \(\mathbf{T}\) 上的一个正函数，关于 \(\mu\) 可测且本性有界，并令 \(\nu = g \cdot \mu\)。证明：若函数 \(\mathbf{f}\) 定义在 \(\mathbf{T}\) 上、取值于一个 Banach 空间或 \(\overline{\mathbf{R}}\) 中，且为 \(\mu\) -可积，则 \(\mathbf{f}\) 为 \(\nu\) -可积。（注意对某个常数有 \(\nu \leqslant a\mu\)，其中 \(a > 0\)，并利用 Ch. IV, §1, No.~3 的命题 15。）
\item
  设 X 为局部紧空间，\(\Lambda : t \mapsto \lambda_t\) 与 \(\Lambda': t \mapsto \lambda_t'\) 为两个 \(\mu\) -适宜映射，把 T 映入 \(\mathcal{M}_+(X)\)，并令 \(\nu = \int \lambda_t d\mu(t)\)，\(\nu' = \int \lambda_t' d\mu(t)\)。证明：若关于 \(\mu\) 局部几乎处处 \(\lambda_t'\) 是以 \(\lambda_t\) 为基的测度，则 \(\nu'\) 是以 \(\nu\) 为基的测度。（考察 X 的一个 \(\nu\) -可忽略紧子集 K。它在局部几乎处处是 \(\lambda_t\) -可忽略的，从而在局部几乎处处是 \(\lambda_t'\) -可忽略的。应用 §3 的命题 5。）
\item
  \begin{enumerate}
  \def\labelenumii{\alph{enumii})}
  \tightlist
  \item
    设 \(\mu\) 为 T 上的一个适度测度 \(\geqslant 0\)。证明存在一个在 T 上连续的函数 \(h \geqslant 0\)，使得 \(h \cdot \mu\) 有界且与 \(\mu\) 等价（仿照命题 11 的方式论证）。
  \end{enumerate}
\end{enumerate}

\begin{enumerate}
\def\labelenumi{\alph{enumi})}
\setcounter{enumi}{1}
\tightlist
\item
  设 \((\mu_{n})\) 为 T 上的一个由适度测度 \(\geqslant0\) 组成的序列。证明存在 T 上的一个测度 \(\mu\geqslant0\)，使每个 \(\mu_{n}\) 都以 \(\mu\) 为基（归结为每个 \(\mu_{n}\) 都有界的情形）。
\end{enumerate}

\begin{enumerate}
\def\labelenumi{\arabic{enumi})}
\setcounter{enumi}{4}
\item
  设 \(\rho, \sigma\) 为 T 上的两个原子测度，并设 M、N 分别为承载 \(|\rho|\) 与 \(|\sigma|\) 的最小集。为使 \(\rho\) 与 \(\sigma\) 互奇异，必须且只需 M \(\cap\) N = \(\varnothing\)。由此导出一个原子测度 \(\nu\) 的例子，它定义在 R 的区间 I = {[}0,1{]} 上，使 I 既是 \(\nu^{+}\) 的支集又是 \(\nu^{-}\) 的支集。
\item
  \begin{enumerate}
  \def\labelenumii{\alph{enumii})}
  \tightlist
  \item
    设 \(\mu\) 为 T 上的正测度。若 \(A \subset T\) 普遍可测且非局部 \(\mu\) -可忽略，证明存在一个由 A 承载的正测度 \(\nu\)，使得 \(\nu \neq 0\) 且 \(\nu \leqslant \mu\)。（取 \(\nu = \varphi_{K} \cdot \mu\)，其中 K 为 A 的一个非 \(\mu\) -可忽略紧子集。）
  \end{enumerate}
\end{enumerate}

\begin{enumerate}
\def\labelenumi{\alph{enumi})}
\setcounter{enumi}{1}
\item
  设 \(\mathbf{M}\) 为 \(\mathbf{T}\) 的一个普遍可测子集。为使正测度 \(\lambda\) 在 \(\mathbf{T}\) 上由 \(\mathbf{M}\) 承载，必须且只需 \(\lambda\) 与每个由 \(\mathbb{C}\mathbf{M}\) 承载的正测度都互奇异（利用 \(a\)）。
\item
  由 \(b\) ) 推出：测度 \(\rho\)（在 \(T\) 上）当其 \(|\rho|\) 由 \(M\) 承载时，构成 \(\mathcal{M}(T)\) 中的一个带。这个带在 \(\mathcal{M}(T)\) 中是淡闭的，若 \(M\) 为 \(T\) 的闭子集（Ch. III, §2, 命题 6）。
\item
  证明 I = {[}0,1{]} 上的 Lebesgue 测度是由一个固定可数集 A ⊂ I 承载的一列原子测度的淡极限（参见 Ch. III, §2, 定理 1 与命题 13）。因此由 A 承载的测度所成的带不是淡闭的。
\end{enumerate}

¶ 7) a) 设 \(\mu\) 为 T 上的正测度，A 为一个 \(\mu\) -可积集且使 \(\mu(A) > 0\)。证明：若对每个 \(\mu\) -可积集 \(B \subset A\)，或者 \(\mu(B) = 0\)，或者 \(\mu(B) = \mu(A)\)，则存在一点 \(a \in A\) 使 \(\mu(\{a\}) = \mu(A)\)。（考察所有紧集 \(K \subset A\)（其中 \(\mu(K) = \mu(A)\)）的交；证明它非空、测度等于 \(\mu(A)\)，且退化为一点。）

\begin{enumerate}
\def\labelenumi{\alph{enumi})}
\setcounter{enumi}{1}
\tightlist
\item
  假设 \(\mu\) 是弥散测度。对每个 \(\mu\) -可积集 A，若 \(\mu(A) > 0\)，又对每个 \(\varepsilon > 0\)，证明存在 A 的一个 \(\mu\) -可积子集 B，使 \(0 < \mu(B) \leqslant \varepsilon\)（借助 \(a\)）注意到存在 A 的一个 \(\mu\) -可积子集 C，使 \(0 < \mu(C) \leqslant \frac{1}{2} \mu(A)\)）。由此推出：当 X 跑遍 A 的 \(\mu\) -可积子集全体时，\(\mu(X)\) 的值集是闭区间 \([0, \mu(A)]\)（对每个数 \(\beta\)，这里 \(0 < \beta < \mu(A)\)，设 \(\gamma\) 为 A 中使 \(\mu(X) \leqslant \beta\) 的可测子集 X 的测度的上确界；先利用前述结果证明 \(\gamma = \beta\)，再证明存在 A 的可测子集的递增序列 \((X_n)\)，使 \(\lim_{n \to \infty} \mu(X_n) = \beta\)）。
\end{enumerate}

¶ 8) a) 设 \(\nu\) 为 T 上的正原子测度，A 为 T 的一个 \(\nu\) -可积子集。证明 \(\nu(X)\)（X 跑遍 A 的 \(\nu\) -可积子集全体）的值集在 R 中闭。（设 M 为承载 \(\nu\) 的最小集。假设 \(A \cap M\) 无穷，并把 \(A \cap M\) 的点排成一个序列 \((a_n)\)，考察映射 \(\varphi\)：它把积空间 \(\{0,1\}^N\) 映入 R，由 \(\varphi((\varepsilon_n)) = \sum_n \varepsilon_n \nu(\{a_n\})\) 定义，并证明

它是连续的。）

\begin{enumerate}
\def\labelenumi{\alph{enumi})}
\setcounter{enumi}{1}
\item
  由 a) 与习题 7 b) 推出：若 \(\mu\) 为 T 上任一正测度，A 为 T 的一个 \(\mu\) -可积子集，则 \(\mu(X)\)（X 跑遍 A 的 \(\mu\) -可积子集全体）的值集在 R 中闭。把这一结果推广到 \(\mu\) 为 T 上任一实测度的情形。
\item
  由习题 7 b) 推出：若 \(\mu\) 为 T 上的弥散实测度，则 \(\mu(X) = \mu^{+}(X) - \mu^{-}(X)\)（X 跑遍 T 的 \(|\mu|\) -可积子集全体）的值集是 R 的一个闭区间（有界或无界）。
\item
  给出一个原子正测度 \(\nu\) 的例子，它在非紧局部紧空间 \(T\) 上，使 \(\nu(X)\) 的值集不闭，这里 \(X\) 跑遍 \(\nu\) -可积子集（\(T\) 中）的全体（取 \(\nu\) 使 \(\inf_{t \in T} \nu(\{t\}) > 0\)）。
\end{enumerate}

\begin{enumerate}
\def\labelenumi{\arabic{enumi})}
\setcounter{enumi}{8}
\item
  设 \(\mu\) 与 \(\nu\) 为 T 上两个互奇异的正测度。证明：对每个数 \(p\)（\(1 \leqslant p \leqslant +\infty\)），拓扑向量空间 \(\mathrm{L}^p(\mathrm{T}, \mu + \nu)\) 同构于拓扑向量空间 \(\mathrm{L}^p(\mathrm{T}, \mu)\) 与 \(\mathrm{L}^p(\mathrm{T}, \nu)\) 的积空间。
\item
  \begin{enumerate}
  \def\labelenumii{\alph{enumii})}
  \tightlist
  \item
    设 \(\mu\) 为一个 \(\neq 0\) 的弥散正测度，定义在局部紧空间 \(T\) 上。证明在 \(\mathcal{L}^1(T, \mu)\) 中存在一个函数序列 \(f_n\)，使 \(N_1(f_n) = 1\) 对一切 \(n\) 成立，且弥散测度序列 \(f_n \cdot \mu\) 淡收敛于一个点测度 \(\varepsilon_a\)。由此推出 \(L^1(T, \mu)\)（从而 \(L^\infty(T, \mu)\) 亦然）不是自反的。
  \end{enumerate}
\end{enumerate}

\begin{enumerate}
\def\labelenumi{\alph{enumi})}
\setcounter{enumi}{1}
\item
  设 \(\nu\) 为 T 上的一个支集无穷的正原子测度。设 A 为承载 \(\nu\) 的最小集，B 为 A 的一个可数无穷子集，\(\lambda\) 为测度 \(\varphi_{\mathrm{B}} \cdot \nu\)。证明 \(\mathrm{L}^{\infty}(\mathrm{T}, \lambda)\) 的对偶不是可分的，从而不能同构于 \(\mathrm{L}^{1}(\mathrm{T}, \lambda)\)（参见 TVS, I, §2, 习题 1）。
\item
  由 a) 与 b) 推出：为使局部紧空间 T 上的正测度 \(\mu\) 使 \(L^{1}(T,\mu)\) 自反，必须且只需 \(\mu\) 的支集有限（利用习题 9）。
\end{enumerate}

¶ 11) 设 S 为紧 Stone 空间（Ch. II, §1, 习题 13 f)）；称 S 上的正测度 \(\mu\) 为正规测度，如果对 S 上连续函数的每个递增有向族 \((f_{\alpha})_{\alpha \in A}\)，它在 \(\mathcal{C}(S)\) 中有上界，且其在完全格序空间 \(\mathcal{C}(S)\) 中的上确界为 \(f\)（该上确界未必等于族 \((f_{\alpha})\) 的上包络），则有 \(\mu(f) = \sup_{\alpha \in A} \mu(f_{\alpha})\)。称 S 上的实测度 \(\lambda\) 为

正规测度，如果正测度 \(\lambda^{+}\) 与 \(\lambda^{-}\) 都是正规的。

\begin{enumerate}
\def\labelenumi{\alph{enumi})}
\item
  为使 S 上的正测度 \(\mu\) 正规，必须且只需每个无处稠密集 A 都是 \(\mu\) -可忽略的（为证条件必要，考察 S 中包含 A 的既开又闭的集；为证条件充分，注意在 \(\mathcal{C}(S)\) 中，一个有上界的递增有向族 \((f_{\alpha})\) 的上包络与其上确界在一个贫集的余集上相等（参见 Ch. II, §1, 习题 13 f)））。由此推出 \(\mu\) 的支集既开又闭。
\item
  设 \(\mu\) 为 S 上的正规正测度，\(f\) 为 \(\mu\) -可测数值函数，\(g\) 为 S 上 \(\leqslant f\) 的最大下半连续函数。证明 \(f\) 与 \(g\) 关于 \(\mu\) 几乎处处相等（参见 Ch. IV, §5, 习题 16 b)）。由此推出
\end{enumerate}

对 S 的每个 \(\mu\) -可测子集 A，A - A 与 \(\overline{A}\) -A 都是 \(\mu\) -可忽略的。

\begin{enumerate}
\def\labelenumi{\alph{enumi})}
\setcounter{enumi}{2}
\tightlist
\item
  证明：在 S 上测度所成的 Banach 空间 \(\mathcal{M}(\mathrm{S})\) 中，正规测度全体是一个闭线性子空间，并且是 \(\mathcal{M}(\mathrm{S})\) 的完全格序空间结构下的一个带。
\end{enumerate}

¶ 12) 称紧 Stone 空间 H 为超 Stone 空间，如果 H 上（习题 11 的）正规正测度的支集的并在 H 中稠密。

\begin{enumerate}
\def\labelenumi{\alph{enumi})}
\item
  设 T 为局部紧空间，\(\mu\) 为 T 上的正测度。证明存在一个同构 \(u \mapsto \theta_{u}\)，它把 Banach 空间 \(\mathrm{L}^{\infty}(\mathrm{T}, \mu)\) 映到某个超 Stone 空间 H 上实值连续函数的空间 \(\mathcal{C}(\mathrm{H})\) 上，且使 \(\theta_{u}^{+} = \theta_{u^{+}}\)（参见命题 14 及 Ch. II, §1, 习题 13）。
\item
  设 H 为紧超 Stone 空间，\((\mu_{\iota})_{\iota\in I}\) 为 H 上的一族正规正测度，使诸 \(\mu_{\iota}\) 的支集的并在 H 中稠密。证明：为使一个集在 H 中无处稠密，必须且只需它为 \(\mu_{\iota}\) -可忽略，对一切 \(\iota\in I\)（注意若 A 是 \(\mu_{\iota}\) -可忽略的，则 \(\overline{A}\) 亦然（习题 11 b)））。由此推出：在超 Stone 空间中，每个贫集都无处稠密。
\item
  在与 \(b\) ) 相同的假设下，证明：若 \(f\) 是一个数值函数，它为 \(\mu_{\iota}\) -可测（对所有 \(\iota \in I\)），则存在一个连续函数 \(g\)（定义在 \(H\) 上），使 \(f\) 与 \(g\) 在一个无处稠密集的余集上相合（利用 \(b\) ) 与习题 11 \(b\)）。
\item
  设 H 为超 Stone 空间。证明存在 H 的一族非空、两两不相交、既开又闭的子集 \((\mathrm{G}_{\alpha})_{\alpha\in\mathrm{A}}\)，其并 T 在 H 中稠密，且对每个 \(\alpha\in A\)，存在一个正规正测度 \(\mu_{\alpha}\)，它定义在 \(G_{\alpha}\) 上并以 \(G_{\alpha}\) 为支集（对 H 上支集两两不相交的正规正测度族所成之集应用 Zorn 引理，并利用习题 11 a)）。设 \(\mu\) 为 T 上在每个 \(G_{\alpha}\) 上的限制为 \(\mu_{\alpha}\) 的测度（Ch. III, §2, 命题 1）。证明：使每个函数 \(f\in\mathcal{C}(\mathrm{H})\) 对应于其在 T 上的限制在 \(\mathrm{L}^{\infty}(\mathrm{T},\mu)\) 中的类的映射，是 Banach 空间 \(\mathcal{C}(\mathrm{H})\) 到 \(\mathrm{L}^{\infty}(\mathrm{T},\mu)\) 上的一个同构（利用 c) 证明该映射是满射）。
\end{enumerate}

¶ 13) 设 T 为局部紧空间，\(\mu\) 为 T 上的正测度，\((f_n)\) 为 \(\mathcal{L}^1(\mathrm{T},\mu)\) 中的函数序列；置 \(\mu_n = f_n \cdot \mu\)，且对每个 \(\mu\) -可测集 A，

\[
\mu_ {n} (\mathrm{A}) = \mu_ {n} ^ {+} (\mathrm{A}) - \mu_ {n} ^ {-} (\mathrm{A}) = \int_ {\mathrm{A}} f _ {n} d \mu .
\]

\begin{enumerate}
\def\labelenumi{\alph{enumi})}
\item
  证明：若 \(\mathbf{T}\) 不是紧的，则可能发生这样的情形：序列 \((f_n)\) 在 \(\mathcal{L}^1 (\mathrm{T},\mu)\) 中不有界，但积分序列 \(\int f_n g d\mu\) 对每个 \(g\in \mathcal{K}(\mathrm{T})\) 有界。
\item
  证明：若序列 \((\mu_n(A))\) 有界——对每个退化为一点的子集 \(A\)（在 \(T\) 中），以及对每个开集 \(A\)（\(\mu\) 在 \(A\) 的边界上诱导的测度具有有限支集）——则序列 \((f_n)\) 在 \(\mathcal{L}^1\) 中有界，或等价地说，范数序列 \((\|\mu_n\|)\) 有界。（先证每一点 \(t_0\)（\(T\) 中的）都具有开邻域 \(U\)，使数列 \(|\mu_n|(U)\) 有界。为此，用反证法证明：在相反情形，可以构造一个严格递增的整数序列 \((n_k)\)、一个递减序列 \((U_k)\)（由 \(t_0\) 的邻域组成），以及开集的一个序列 \((W_k)\)（对 \(\mu\) 可求积，Ch. IV, §5, 习题 17），它们具有下列性质：\(\overline{U}_k \subset U_{k-1}\)，\(\mu(U_k - \{t_0\}) \leqslant 1/k\)，\(|\mu_{n_i}|(U_k - \{t_0\}) \leqslant 1\)（对 \(i < k\)），\(\overline{W}_k \subset U_k - \overline{U}_{k+1}\)，最后
\end{enumerate}

\[
\left| \mu_ {n _ {k}} \left(\mathrm{W} _ {k}\right) \right| > k + \sum_ {i <   k} \left| \mu_ {n _ {k}} \left(\mathrm{W} _ {i}\right) \right|.
\]

最后，考察诸 \(W_{k}\) 的并 W，以得出矛盾（“滑峰法”）。然后以类似的论证证明：存在 T 的一个紧子集 K，使序列 \(|\mu_{n}|(T - K)\) 有界。）

\begin{enumerate}
\def\labelenumi{\alph{enumi})}
\setcounter{enumi}{2}
\item
  由 \(b\) ) 推出：若对每个开集 \(A\)（在 \(T\) 中），序列 \((\mu_n(A))\) 都有界，则序列 \((f_n)\) 在 \(\mathcal{L}^1\) 中有界。
\item
  在区间 \([0,1]\) 上，若取 \(\mu_{n}\) 为由点 t=0 处的质量 n 与点 \(t=1/n\) 处的质量 -n 所定义的测度，则序列 \((\mu_{n}(A))\) 对每个关于所有测度 \(|\mu_{n}|\) 都可求积的开集 A 有界；然而范数序列 \(\|\mu_{n}\|\) 不有界。类似地，若取 \(\mu_{n}\) 为由点 \(t=1/n\) 处的质量 n 与点 \(t=1/(n+1)\) 处的质量 -n 所定义的测度，则序列 \((\mu_{n}(A))\) 对每个含于 \([0,1]\) 的区间 A（从而对每个有限子集 A）有界，而序列 \((\|\mu_{n}\|)\) 不有界。
\end{enumerate}

\begin{enumerate}
\def\labelenumi{\arabic{enumi})}
\setcounter{enumi}{13}
\tightlist
\item
  设 \(\theta\) 为空间 \(\mathcal{L}^{\infty}(\mathrm{T},\mu)\) 上的一个线性形式；为使 \(\theta\) 具有 \(f\mapsto\int fg d\mu\) 的形式（其中 \(g\in\mathcal{L}^{1}(\mathrm{T},\mu)\)），必须且只需 \(\theta\) 满足下列条件：\(1^{\circ}\) 对每个 \(\varepsilon>0\)，存在 \(\delta>0\)，使关系 \(\mu^{*}(\mathrm{A})\leqslant\delta\) 蕴含 \(|\theta(h)|\leqslant\varepsilon\)，对每个 \(\mu\) -可测函数 \(h\)（\(|h|\leqslant\varphi_{\mathrm{A}}\)）；\(2^{\circ}\) 对每个 \(\varepsilon>0\)，存在紧集 K，使对每个 \(\mu\) -可积集 B \(\subset\) T−K，有 \(|\theta(\varphi_{\mathrm{B}})|\leqslant\varepsilon\)。（利用 Lebesgue--Nikodym 定理。）
\end{enumerate}

¶ 15) 设 H 为空间 \(\mathcal{L}^1 (\mathrm{T},\mu)\) 的一个有界子集，\(\widetilde{\mathrm{H}}\) 为它在 Banach 空间 \(\mathrm{L}^1 (\mathrm{T},\mu)\) 中的典范像。证明下列条件等价：

\(\alpha)\) 下列两个条件成立：

\(\alpha_{1})\) 对每个 \(\varepsilon > 0\)，存在紧集 \(L \subset T\)，使 \(\int_{T - L} |f| d\mu \leqslant \varepsilon\) 对一切 \(f \in H\) 成立；

\(\alpha_{2})\) 对每个 \(\varepsilon > 0\)，存在 \(\eta > 0\)，使对 T 的每个满足 \(\mu^{*}(\mathrm{U}) \leqslant \eta\) 的开集 U，有 \(\int_{\mathrm{U}} |f| d\mu \leqslant \varepsilon\)，对一切 \(f \in \mathrm{H}\)。

\(\beta)\) 条件 \(\alpha_{1}\) ) 与下列条件 \(\beta_{2}\) ) 都成立：

\(\beta_{2})\) 对每个紧集 \(\mathbf{K} \subset \mathbf{T}\) 及每个 \(\varepsilon > 0\)，存在开邻域 \(\mathbf{U}\)（含 \(\mathbf{K}\)），使 \(\int_{\mathbf{U} - \mathbf{K}} |f| d\mu \leqslant \varepsilon\) 对一切 \(f \in \mathbf{H}\) 成立。

\(\gamma)\) 对每个函数序列 \((g_n)\)（属于 \(\mathcal{L}^{\infty}\)，一致有界且依测度收敛于一个函数 \(g\)），有 \(\lim_{n\to \infty}\int fg_n d\mu = \int fg d\mu\)，对 \(f\) 跑遍 H 一致地成立。

\(\delta)\) 对每个由 T 上连续且在无穷远点趋于 0 的函数组成的一致有界序列 \((h_n)\)，若 \(\lim_{n\to \infty}h_n(t) = 0\) 对每个 \(t\in \mathbb{T}\) 成立，则 \(\lim_{n\to \infty}\int fh_n d\mu = 0\) 对 \(f\) 跑遍 H 一致地成立。

\(\zeta)\) 对 T 的两两不相交开集的每个无穷序列 \((\mathbf{U}_n)\)，有 \(\lim_{n\to \infty}\int_{\mathbf{U}_n}f d\mu = 0\)，对 \(f\) 跑遍 H 一致地成立。

（为证 \(\zeta\) ) 蕴含 \(\beta\) ) 且 \(\beta\) ) 蕴含 \(\alpha\) )，按习题 13 \(b\) ) 的方式使用“滑峰法”。然后证明 \(\alpha\) ) 蕴含 \(\gamma\) )，\(\gamma\) 蕴含 \(\delta\) )，以及 \(\delta\) ) 蕴含 \(\zeta\)。）

若再假设：对每个 \(t \in \mathbf{T}\)（满足 \(\mu(\{t\}) \neq 0\)），\(f(t)\) 所成的集在 \(\mathbf{R}\) 中有界（当 \(f\) 跑遍 \(\mathbf{H}\)），试证上述条件与下列条件等价：

\(\theta)\) H 是 \(\mathbf{L}^1\) 关于弱化拓扑 \(\sigma (\mathbf{L}^{1},\mathbf{L}^{\infty})\) 的一个相对紧子集

（为证 \(\theta\) ) 蕴含 \(\beta\) )，用反证法，并利用 Šmulian 定理（TVS, IV, §5, No.~3, 定理 2）与“滑峰法”。为证 \(\alpha\) ) 蕴含 \(\theta\) )，先证 H 在 \(\mathcal{L}^1\) 中有界，然后应用 Eberlein 定理（TVS, IV, §5, No.~3, 定理 1），并利用习题 14。）

¶ 16) a) 设 \((f_n)\) 为 \(\mathcal{L}^1(\mathrm{T},\mu)\) 中函数的一个序列。证明下列条件等价：

\(\alpha)\) 序列 \((f_n)\) 在 \(\mathbf{L}^1\) 中关于弱化拓扑 \(\sigma(\mathbf{L}^1, \mathbf{L}^\infty)\) 收敛

\(\beta)\) \(f_{n}\) 所成之集在 \(\mathbf{L}^{1}\) 中关于弱化拓扑相对紧，且测度序列 \(f_{n} \cdot \mu\) 在 \(\mathcal{M}(\mathrm{T})\) 中淡收敛。

\(\gamma)\) 对 T 的每个开子集 U，数列 \(\int_{\mathrm{U}} f_n d\mu\) 在 \(\mathbf{R}\) 中收敛。

（为证 \(\beta\) ) 蕴含 \(\alpha\) )，利用习题 15 的判别准则 \(\alpha\) ) 与可测函数的定义。为证 \(\gamma\) ) 蕴含 \(\beta\) )，先考察空间 \(\mathrm{L}^1(\mathbf{N})\)（对每点由质量 \(+1\) 定义的离散测度）的特殊情形，利用 TVS, IV, §2 的习题 15。在一般情形，应用习题 15 的判别准则 \(\zeta\) ) 以化归为 \(\mathrm{L}^1(\mathbf{N})\) 的情形，办法是对每个 \(f_n\) 联系可和序列 ( \(\int_{\mathrm{U}_m} f_n d\mu)_{m \in \mathbf{N}}\)。）

\begin{enumerate}
\def\labelenumi{\alph{enumi})}
\setcounter{enumi}{1}
\item
  由这些结果推出：在 \(L^{1}\) 中，每个关于弱化拓扑的 Cauchy 序列对该拓扑收敛。
\item
  在区间 T = {}0,1{} 上，设 \(\mu\) 为 Lebesgue 测度，且对每个整数 \(n \geqslant 1\)，设 \(\mu_{n}\) 为在每点 k/n（\(0 \leqslant k < n\)）由质量 1/n 定义的测度；设 \(\nu\) 为 T 上的正测度，使 \(\mu = g \cdot \nu\)，\(\mu_{n} = f_{n} \cdot \nu\)（习题 4 b)）。证明序列 \((\mu_{n})\) 淡收敛于 \(\mu\)，对 T 的每个有限子集 A 以及每个使 \(\mu_{n}(A)\) 在其边界上诱导的测度具有有限支集的开集 A，\(\mu(A)\) 趋于 \(\nu\)，但对所有开集 \(\mu_{n}(U)\)，\(\mu(U)\) 并不都趋于 \(U \subset T\)（参见习题 13 b)）。
\end{enumerate}

\begin{enumerate}
\def\labelenumi{\arabic{enumi})}
\setcounter{enumi}{16}
\tightlist
\item
  \begin{enumerate}
  \def\labelenumii{\alph{enumii})}
  \tightlist
  \item
    设 \(\mu\) 为区间 \(T = [0, +\infty[\) 上的 Lebesgue 测度。证明：若 \(1 \leqslant r < s < +\infty\)，则 \(L^r \cap L^s\) 上由 \(L^r\) 与 \(L^s\) 的弱化拓扑诱导的拓扑不可比较（参见 Ch. IV, §6, 习题 8）。若 \(f_n\) 为区间 \([n, n+1]\) 的特征函数，证明序列 \((\widetilde{f}_n)\) 对一切 \(L to the p\)（\(p > 1\)）的弱化拓扑趋于 0，但对 \(L^1\) 的弱化拓扑不然。
  \end{enumerate}
\end{enumerate}

\begin{enumerate}
\def\labelenumi{\alph{enumi})}
\setcounter{enumi}{1}
\item
  设 \(\lambda\) 为区间 \(T = [0,1]\) 上的 Lebesgue 测度。证明：若 \(r < s\)，则 \(L^s\) 的弱化拓扑严格细于 \(L^s\) 上由 \(L^r\) 的弱化拓扑诱导的拓扑（参见 Ch. IV, §6, 习题 8）。
\item
  设 \(\mu\) 为离散空间上的正测度，每点的测度均为 1。证明：若 \(1 \leqslant r < s < +\infty\)，则 \(\mathbf{L}^r\) 的弱化拓扑严格细于 \(\mathbf{L}^r\) 上由 \(\mathbf{L}^s\) 的弱化拓扑诱导的拓扑（Ch. IV, §6, 习题 9）。
\end{enumerate}

\begin{enumerate}
\def\labelenumi{\arabic{enumi})}
\setcounter{enumi}{17}
\tightlist
\item
  设 A 为含于 \(L^{r} \cap L^{s}\) 的一个集，它在 \(L^{r}\) 与 \(L^{s}\) 中都有界（\(1 < r < s < +\infty\)）。证明：对每个使 \(r \leqslant p \leqslant s\) 的 p，A 在 \(L^{p}\) 中有界，且诸 \(L^{p}\) 的弱化拓扑在 A 上诱导的拓扑相同（注意 \(\widetilde{\mathcal{K}}(\mathrm{T})\) 在每个空间 \(L^{q}\)（\(1 \leqslant q < +\infty\)）中都稠密）。证明：若 r = 1，该性质不再成立（参见习题 17 a)）。
\end{enumerate}

¶ 19) 设 \(\mu\) 为 T 上的正测度。

\begin{enumerate}
\def\labelenumi{\alph{enumi})}
\item
  设 \((f_n)\) 为 \(\mathcal{L}^p\)（\(1 < p < +\infty\)）中函数的一个序列，在 \(\mathcal{L}^p\) 中有界且依测度收敛于 0。证明序列 \((\widetilde{f}_n)\) 在 \(\mathbf{L}^p\) 中关于弱化拓扑收敛于 0。
\item
  给出函数序列 \((f_n)\) 的一个例子：它在 \(\mathcal{L}^1\) 中、在 \(\mathcal{L}^1\) 中有界且依测度收敛于 0，但序列 \((\widetilde{f}_n)\) 对 \(\mathrm{L}^1\) 的弱化拓扑不收敛于 0（取 \(\mu\) 为 \(\mathrm{T} = [0,1]\) 上的 Lebesgue 测度，并利用习题 15 的判别准则 \(\alpha\) )）。
\item
  设 \(\mu\) 为 \(T = [0,1]\) 上的 Lebesgue 测度。证明：若 \(f_{n}(t) = \sin nt\)，则序列 \((\widetilde{f}_n)\) 对一切 \(L^{p}\)（\(1 \leqslant p < +\infty\)）的弱化拓扑收敛于 0，但 \(f_{n}\) 不依测度收敛于 0。
\end{enumerate}

\begin{enumerate}
\def\labelenumi{\arabic{enumi})}
\setcounter{enumi}{19}
\tightlist
\item
  \begin{enumerate}
  \def\labelenumii{\alph{enumii})}
  \tightlist
  \item
    设 \((f_n)\) 为 \(\mathcal{L}^1\) 中函数的一个序列，满足：\(1^\circ\) 序列 \((f_n)\) 依测度收敛于 \(0; 2^\circ\) 对每个 \(\varepsilon > 0\)，存在 \(\eta > 0\)，使关系 \(\mu^*(A) \leqslant \eta\) 蕴含 \(\limsup_{n \to \infty} \int_A^* |f_n| d\mu \leqslant \varepsilon; 3^\circ\) 对每个 \(\varepsilon > 0\)，存在一个可积集 \(B\)，使 \(\int_{T - B} |f_n| d\mu \leqslant \varepsilon\) 对一切 \(n\) 成立。证明序列 \((f_n)\) 平均收敛于 \(0\)。
  \end{enumerate}
\end{enumerate}

\begin{enumerate}
\def\labelenumi{\alph{enumi})}
\setcounter{enumi}{1}
\tightlist
\item
  设 \(\mu\) 为 \(T = [0,1]\) 上的 Lebesgue 测度。对 \(1 < p < +\infty\)，用 \(f_n\) 表示等于 \(n^{2/p}\)（在区间 \(\left[\frac{1}{n+1}, \frac{1}{n}\right]\) 内）而在其外等于 0 的函数。证明序列 \((f_n)\) 满足 \(a)\) 的三个条件，且序列 \((\widetilde{f}_n)\) 对 \(L to the p\) 的弱化拓扑收敛于 0，但对 \(p\) 阶平均收敛拓扑不然。
\end{enumerate}

¶ 21) a) 设 \((f_n)\) 为 \(\mathcal{L}^1\) 中函数的一个序列，满足：\(1^\circ \liminf_{n \to \infty} f_n(t) \geqslant 0\) 在 T 中几乎处处成立；\(2^\circ\) 对每个 \(\varepsilon > 0\)，存在 \(\eta > 0\)，使关系 \(\mu^*(A) \leqslant \eta\) 蕴含 \(\limsup_{n \to \infty} \int_A^* |f_n| d\mu \leqslant \varepsilon\)；\(3^\circ\) 对每个 \(\varepsilon > 0\)，存在一个可积集 B，使 \(\int_{T-B} |f_n| d\mu \leqslant \varepsilon\) 对一切 \(n\) 成立；\(4^\circ \lim_{n \to \infty} \int f_n d\mu = 0\)。证明序列 \((f_n)\) 平均收敛于 0。（注意，对每个 \(\varepsilon > 0\)，若 \(C_m\) 为 \(t \in T\) 中使 \(f_n(t) \geqslant -\varepsilon\)（对至少一个 \(n \geqslant m\)）成立的点所成之集，则序列 \((C_m)\) 递减且其交可忽略。）

\begin{enumerate}
\def\labelenumi{\alph{enumi})}
\setcounter{enumi}{1}
\tightlist
\item
  设 \((f_n)\) 为 \(\mathcal{L}^1\) 中函数的一个序列，满足：\(1^\circ\) 存在一个可积函数 \(g\)，使 \(f_n \geqslant g\) 对一切 \(n\) 成立；\(2^\circ \liminf_{n \to \infty} f_n(t) \geqslant 0\) 几乎处处成立；\(3^\circ \limsup_{n \to \infty} \int f_n d\mu \leqslant 0\)。证明序列 \((f_n)\) 平均收敛于 0。（先证对每个可测集 A，\(\liminf_{n \to \infty} \int_{A} f_n d\mu \geqslant 0\)；然后注意，对每个可测集 A，
\end{enumerate}

\[
\limsup _ {n \to \infty} \int f _ {n} d \mu \geqslant \limsup _ {n \to \infty} \int_ {\mathrm{A}} f _ {n} d \mu + \liminf _ {n \to \infty} \int_ {\mathrm{T-A}} f _ {n} d \mu .
\]

¶ 22) 在 Banach 空间 \(\mathrm{E} = \mathcal{M}^1(\mathrm{T})\)（局部紧空间 \(\mathrm{T}\) 上有界测度所成者）中，设 \(\mathrm{H}\) 为关于弱化拓扑的一个紧集。

\begin{enumerate}
\def\labelenumi{\alph{enumi})}
\item
  证明：对每个 \(\varepsilon > 0\)，存在 T 上的一个正测度 \(\mu_{\varepsilon}\)，使每个测度 \(\nu \in H\) 都是一个以 \(\mu_{\varepsilon}\) 为基的正测度与一个测度 \(\lambda\)（它与 \(\mu_{\varepsilon}\) 互奇异且范数 \(\leqslant \varepsilon\)）之和。（利用如下事实：空间 E 赋以其范数与序结构后同构于一个空间 \(L^{1}(S, \rho)\)，其中 S 是紧 Stone 空间的并，而 \(\rho\) 是 S 上的正测度（Ch. IV, §4, 习题 10），并在 \(L^{1}(S, \rho)\) 中应用习题 15 的判别准则 \(\alpha_{1}\) ）。）
\item
  由 a) 推出：存在 T 上的一个正测度 \(\mu\)，使所有测度 \(\nu \in H\) 都以 \(\mu\) 为基（利用习题 4 b)）。
\end{enumerate}

\begin{enumerate}
\def\labelenumi{\arabic{enumi})}
\setcounter{enumi}{22}
\item
  设 \(\mu\) 为局部紧空间 \(T\) 上的正测度，\(p\) 为使 \(1 \leqslant p < +\infty\) 的数，\(q\) 为共轭指数。证明：若 \(g\) 是有限可测函数，使对每个函数 \(f \in \mathcal{L}^p\)，函数 \(fg\) 可积，则 \(g\) 局部几乎处处等于 \(\mathcal{L}^q\) 中的一个函数。（利用闭图定理（TVS, I, §3, 定理 1 的推论 5）证明映射 \(f \mapsto fg\)（它把 \(\mathcal{L}^p\) 映入 \(\mathcal{L}^1\)）连续。）
\item
  设 \(p, q\) 为两个共轭指数。若 B（相应地，C）是 \(L to the p\)（相应地，\(L^q\)）的有界子集，证明 \(B \times C\) 到 \(L^1\) 中的映射——它由 \((f, g) \mapsto fg\) 通过过渡到商得到——当 \(L to the p\)、\(L^q\) 与 \(L^1\) 分别赋以拓扑 \(\sigma(L to the p, L^q)\)、\(\sigma(L^q, L to the p)\) 与 \(\sigma(L^1, L^\infty)\) 时，未必连续（参见习题 10）。
\end{enumerate}

¶ 25) a) 若 u 与 v 为任意两个实数，试对 \(1 < p < +\infty\) 证明不等式

\[
| u + v | ^ {p} \leqslant | u | ^ {p} + p v u | u | ^ {p - 2} + a \sum_ {r = 2} ^ {[ p ]} | v | ^ {r} | u | ^ {p - r} + b | v | ^ {p},
\]

其中 \(a\) 与 \(b\) 是两个仅依赖于 \(p\) 的常数，而 \([p]\) 是 \(p\) 的整数部分（FRV, III, §2, 习题 6）。

\begin{enumerate}
\def\labelenumi{\alph{enumi})}
\setcounter{enumi}{1}
\tightlist
\item
  设 \((f_n)\) 为在 \(\mathbf{L}^p\) 中关于弱化拓扑 \((1 < p < +\infty)\) 收敛于 0 的一个序列。证明存在子序列 \((f_{n_k})\)（取自 \((f_n)\)），使在令 \(s_m = \sum_{k=1}^{m} f_{n_k}\) 时，有
\end{enumerate}

\[
\left| \int s _ {m - 1} | s _ {m - 1} | ^ {p - 2} f _ {n _ {m}} d \mu \right| \leqslant 1
\]

（递归地定义序列 \((n_k)\)）。由此，利用 \(a\) 与 Hölder 不等式推出：若 \(p > 2\)，则存在两个常数 \(a, b\) 使

\[
\mathrm{N} _ {1} (| s _ {n} | ^ {p}) \leqslant \mathrm{N} _ {1} (| s _ {n - 1} | ^ {p}) + a + b \mathrm{N} _ {1} (| s _ {n - 1} | ^ {p - 2}),
\]

而若 \(1 < p \leqslant 2\)，则存在一个常数 \(c\) 使

\[
\mathrm{N} _ {1} \big (| s _ {n} | ^ {p} \big) \leqslant \mathrm{N} _ {1} \big (| s _ {n - 1} | ^ {p} \big) + c.
\]

由此得出结论：

\[
\begin{array}{l l} \mathrm{N} _ {p} (s _ {n}) = O (n ^ {1 / 2}) & \text { for } p > 2 \\ \mathrm{N} _ {p} (s _ {n}) = O (n ^ {1 / p}) & \text { for } 1 <   p \leqslant 2. \end{array}
\]

\begin{enumerate}
\def\labelenumi{\alph{enumi})}
\setcounter{enumi}{2}
\tightlist
\item
  证明前面的结果一般不能改进（参见习题 19 与 20 b)）。
\end{enumerate}

¶ 26) 设 \(\mu_1, \ldots, \mu_m\) 为局部紧空间 \(T\) 上有限个测度，它们可写成形式 \(\mu_k = f_k \cdot \mu\)，其中 \(\mu\) 是一个测度 \(\geqslant 0\) 且 \(|f_k| \leqslant 1\)（No.~9）。

\begin{enumerate}
\def\labelenumi{\alph{enumi})}
\item
  设 \(\Phi\) 为定义在 \(\mathbf{R}^m\) 上且正齐次的函数所成之集，且使 \(u(f_1, \ldots, f_m)\) 是局部 \(\mu\) -可积的（对每个函数 \(u \in \Phi\)）。证明：若函数序列 \((u_n)\)（属于 \(\Phi\)）在 \(\mathbf{R}^m\) 的每个紧子集上一致收敛于一个函数 \(u\)，则 \(u(\mu_1, \ldots, \mu_m)\) 是测度序列 \(u_n(\mu_1, \ldots, \mu_m)\) 在 \(\mathcal{M}(\mathrm{T})\) 的拟强拓扑下的极限（Ch. III, §1, 习题 8）。
\item
  假设 \(u \in \Phi\) 在 \(\mathbf{R}^m\) 上连续。证明映射
\end{enumerate}

\[
(\mu_ {1}, \dots , \mu_ {m}) \mapsto u (\mu_ {1}, \dots , \mu_ {m})
\]

把 \((\mathcal{M}(\mathrm{T}))^m\) 映入 \(\mathcal{M}(\mathrm{T})\) 的这一映射对拟强拓扑连续（先考察 \(u\) 为 Lipschitz 函数的情形；然后利用 \(a\) ) 与 Stone--Weierstrass 定理）。

\begin{enumerate}
\def\labelenumi{\alph{enumi})}
\setcounter{enumi}{2}
\tightlist
\item
  假设 \(u\) 在 \(\mathbf{R}^m\) 上连续。设 \(g\) 为 \(\mathcal{H}(\mathrm{T})\) 中任一函数，\(K\) 为其支集。对每个 \(\varepsilon > 0\)，证明存在一个由相对紧集组成的有限开覆盖 \((A_i)\)（它覆盖 \(K\)），使得对每个有限开覆盖 \((B_j)\)（它覆盖 \(K\) 且比 \((A_i)\) 更细），以及对连续映射 \((h_j)\)（从 \(T\) 到 {}0,1{}）的每个族——其中 \(h_j\) 的支集含于 \(B_j\)，且 \(\sum_{j} h_j(t) = 1\) 在 \(K\) 上成立——都有
\end{enumerate}

\[
\left| \int g \cdot u (f _ {1}, \dots , f _ {m}) d \mu - \sum_ {j} u \left(\int g h _ {j} f _ {1} d \mu , \dots , \int g h _ {j} f _ {m} d \mu\right) \right| \leqslant \varepsilon
\]

（先考察 \(u\) 为 Lipschitz 函数且诸 \(f_k\) 连续的情形，然后过渡到 \(u\) 为 Lipschitz 函数且诸 \(f_k\) 局部可积的情形，最后利用 Stone--Weierstrass 定理处理一般情形）。

\begin{enumerate}
\def\labelenumi{\alph{enumi})}
\setcounter{enumi}{3}
\tightlist
\item
  在 \(\mathbf{R}^2\) 中，设 \(u(x_1, x_2)\) 为正齐次函数，它等于 \(\sqrt{x_1^2 + x_2^2}\)（当 \(x_1 / x_2\) 为无理数时），在相反情形等于 0。设 \(\mu\) 为 \(\mathrm{T} = [0, 1]\) 上的 Lebesgue 测度，并设 \(\nu\) 为 \(\mathrm{T}\) 上由 \(d\nu(t) = t d\mu(t)\) 定义的测度。证明：若 \(g\) 是一个连续且 \(\geqslant 0\) 的函数（在 \(\mathrm{T}\) 上），则 \(\sum_{j} u\left( \int gh_j d\mu, \int gh_j d\nu \right)\) 不趋于任何
\end{enumerate}

关于 g 的支集的有限开覆盖所成的有向集的极限。

\begin{enumerate}
\def\labelenumi{\arabic{enumi})}
\setcounter{enumi}{26}
\tightlist
\item
  设 \(\mathbf{X}\) 为局部紧空间，\(\Lambda : t \mapsto \lambda_t\) 为一个标量本质 \(\mu\) -可积映射，把 \(\mathbf{T}\) 映入 \(\mathcal{M}_+(\mathbf{X})\) 中。假设 \(\Lambda\) 对每个测度 \(\varphi_{\mathbf{K}} \cdot \mu\)（其中 \(\mathbf{K}\) 为紧集）都是预适宜的。
\end{enumerate}

\begin{enumerate}
\def\labelenumi{\alph{enumi})}
\item
  证明 \(\Lambda\) 对每个测度 \(\varphi_{\mathrm{A}}\cdot \mu\) 都是预适宜的，其中 \(\mathbf{A}\) 是一个 \(\mu\) -可测子集（含于 \(\mathbf{T}\) 中）（考察 \(\mathbf{A}\) 的紧子集）。
\item
  证明 \(\Lambda\) 对每个测度 \(f \cdot \mu\) 都是预适宜的，其中 \(f\) 为正的、\(\mu\) -可测且只取有限多个值的函数。
\item
  证明 \(\Lambda\) 对每个测度 \(f \cdot \mu\) 都是预适宜的，其中 \(f\) 为 \(\mu\) -可测且介于 0 与 1 之间的函数。由此推出 \(\Lambda\) 是 \(\mu\) -适宜的。
\end{enumerate}

\begin{enumerate}
\def\labelenumi{\arabic{enumi})}
\setcounter{enumi}{27}
\tightlist
\item
  设 \(\mathbf{X}\) 为局部紧空间，\(\Lambda : t \mapsto \lambda_t\) 为一个 \(\mu\) -适宜映射，把 \(\mathbf{T}\) 映入 \(\mathcal{M}_+(\mathbf{X})\) 中。设 \(\mu' = g \cdot \mu\) 为以 \(\mu\) 为基的正测度，且使 \(\Lambda\) 是标量本质 \(\mu'\) -可积的。
\end{enumerate}

\begin{enumerate}
\def\labelenumi{\alph{enumi})}
\item
  证明 \(\Lambda\) 是 \(\mu'\) -可积的（利用习题 27 b)，或利用关系式 \(\mu' = \sup_n \left( \inf(\mu', n\mu) \right)\) 与 §3 的习题 9）。
\item
  证明 \(\nu' = \int \lambda_t d\mu'(t)\) 是以 \(\nu = \int \lambda_t d\mu(t)\) 为基的测度。
\item
  证明映射 \(t \mapsto g(t)\lambda_t\) 是 \(\mu\) -适宜的，且其积分为 \(\nu'\)。
\end{enumerate}

\begin{enumerate}
\def\labelenumi{\arabic{enumi})}
\setcounter{enumi}{28}
\tightlist
\item
  设 X 为局部紧空间。设 \(\Lambda: t \mapsto \lambda_{t}\) 为一个 \(\mu\) -适宜映射，把 T 映入 \(\mathcal{M}_{+}(X)\) 中。设 g 为 X 上的一个函数，为正、普遍可测且局部有界。证明族 \(t \mapsto g \cdot \lambda_{t}\) 是 \(\mu\) -适宜的，且
\end{enumerate}

\[
\int (g \cdot \lambda_ {t}) d \mu (t) = g \cdot \int \lambda_ {t} d \mu (t)
\]

在下列每个条件下成立：a) \(g\) 对测度 \(\int \lambda_t d\mu(t)\) 是适度的；b) 诸测度 \(\lambda_t\) 有界。

\begin{enumerate}
\def\labelenumi{\arabic{enumi})}
\setcounter{enumi}{29}
\tightlist
\item
  \begin{enumerate}
  \def\labelenumii{\alph{enumii})}
  \tightlist
  \item
    设 M 为 C 的有界子集。证明：M 的闭凸包等于包含 M 的闭圆盘的交集。
  \end{enumerate}
\end{enumerate}

\begin{enumerate}
\def\labelenumi{\alph{enumi})}
\setcounter{enumi}{1}
\tightlist
\item
  设 X 为局部紧空间，\(\mu\) 为 X 上使 \(\|\mu\| = \mu(1) = 1\) 的复测度，f 为 X 上的有界连续复函数。证明 \(\mu(f)\) 属于 C 中 \(f(X)\) 的闭凸包络。（利用 a)。）由此给出命题 9 的一个新证明。
\end{enumerate}

\begin{enumerate}
\def\labelenumi{\arabic{enumi})}
\setcounter{enumi}{30}
\tightlist
\item
  设 \(\theta\) 为 \(T\) 上的复测度，\(F\) 为一个 Banach 空间。
\end{enumerate}

\begin{enumerate}
\def\labelenumi{\alph{enumi})}
\tightlist
\item
  证明空间 \(\mathrm{L}_{\mathrm{loc}}^{1}(\mathrm{T},\theta ;\mathrm{F})\) 是完备的。（把 \(|\theta |\) 写成形式 \(\sum_{\alpha}\mu_{\alpha}\)，其中 \((\mu_{\alpha})\) 是 T 上的测度 \(\geqslant 0\) 的可和族，其支集构成两两不相交紧集的一个局部可数族。这定义了 \(\mathrm{L}_{\mathrm{loc}}^{1}(\mathrm{T},\theta ;\mathrm{F})\) 到 \(\prod_{\alpha}\mathrm{L}_{\mathrm{F}}^{1}(\mathrm{T},\mu_{\alpha})\) 中的一个连续映射。设 \(\mathfrak{F}\) 为 \(\mathrm{L}_{\mathrm{loc}}^{1}(\mathrm{T},\theta ;\mathrm{F})\) 上的一个 Cauchy 滤子。它的像
\end{enumerate}

在 \(\prod_{\alpha} \mathrm{L}_{\mathrm{F}}^{1}(\mathrm{T}, \mu_{\alpha})\) 中收敛于一个元素 \((\mathbf{f}_{\alpha})\)。证明诸 \(\mathbf{f}_{\alpha}\) 定义一个 \(\theta\) -可测的

函数 \(\mathbf{f}\)（在 \(\mathrm{T}\) 上），然后证明 \(\mathbf{f}\) 是局部 \(\theta\) -可积的，且 \(\dot{\mathbf{f}}\) 是 \(\mathfrak{F}\) 在 \(\mathrm{L}_{\mathrm{loc}}^{1}(\mathrm{T},\theta ;\mathrm{F}).\) 中的极限。

\begin{enumerate}
\def\labelenumi{\alph{enumi})}
\setcounter{enumi}{1}
\tightlist
\item
  设 p 为实数 \textgreater{} 1。用 \(\mathcal{L}_{\mathrm{loc}}^{p}(\mathrm{T},\theta;\mathrm{F})\) 表示 T 上取值于 F 的 \(\theta\) -可测函数 f 所成的向量空间，其中 \(|f|^{p}\) 是局部 \(\theta\) -可积的。给它赋以半范数 \(\mathbf{f} \mapsto (\int |\mathbf{f}\varphi_{\mathrm{K}}|^{p}d|\theta|)^{1/p}\)，这里 K 跑遍 T 的紧集所成之集。设 \(\mathrm{L}_{\mathrm{loc}}^{p}(\mathrm{T},\theta;\mathrm{F})\) 为相联的 Hausdorff 空间。证明这个空间是完备的。（方法同 a)。）
\end{enumerate}

\setcounter{exercisegroup}{6}
\edef\CurrentExerciseGroup{\arabic{exercisegroup}}
